\chapter{Complexity — Implicit Poly}
%%% generated by the blueprint pipeline from TCSlib.Complexity.SpaceComplexity.ImplicitPoly
\section{Overview}

Declarations of the Lean module \texttt{TCSlib.Complexity.SpaceComplexity.ImplicitPoly}.

\section{Declarations}

\begin{lemma}[Membership in an Indexed Language]
\lean{Complexity.pairEncode_mem_indexLang}
\label{Complexity.pairEncode_mem_indexLang}
\leanok
\uses{Complexity.LogProg.bits_injective, Complexity.indexLang, Turing.pairEncode, Turing.pairEncode_injective}
\difficulty{4}
Let $p$ be any predicate on pairs consisting of a binary string and a natural number.
For every binary string $x$ and every $i\in\mathbb{N}$, the encoding of $x$ paired with
the binary expansion of $i$ belongs to the language of encodings $\langle
x',\operatorname{bits}(i')\rangle$ for which $p(x',i')$ holds if and only if $p(x,i)$
holds.
\end{lemma}

\begin{definition}[Implicit Enumeration States]
\lean{Complexity.ImplicitEnum.St}
\label{Complexity.ImplicitEnum.St}
\leanok
Introduce a finite set of seven states: ask, bit, emit true, emit false, increment the
counter, increment the bit index, and done.
\end{definition}

\begin{definition}[Implicit enumeration machine]
\lean{Complexity.ImplicitEnum.tm}
\label{Complexity.ImplicitEnum.tm}
\leanok
\uses{Complexity.ImplicitEnum.St, Complexity.LogProg.incBAct, Complexity.LogProg.incCAct, Turing.MultiTapeTM}
The machine has one Boolean work tape and states `ask`, `bit`, `emitT`, `emitF`, `incC`,
`incB`, and `done`, and starts in `ask`. In states `emitT` and `emitF` it outputs true
and false, respectively, and moves to `incC`; in states `incC` and `incB` its
transitions are given by the increment actions used in its definition, and in every
other case it halts without output.
\end{definition}

\begin{definition}[Implicit enumeration program]
\lean{Complexity.ImplicitEnum.prog}
\label{Complexity.ImplicitEnum.prog}
\leanok
\uses{Complexity.ImplicitEnum.St, Complexity.ImplicitEnum.tm, Complexity.LogProg.RProg}
The program has one work tape, two deciders, and the states of the implicit enumeration
machine. Its ordinary machine is the specified Boolean one-tape machine; at state `ask`
it calls decider 0 on the whole virtual input with register 0 as its argument,
proceeding to `bit` on a positive answer and `done` on a negative answer. At state `bit`
it calls decider 1 in the same mode with the same register argument, proceeding to
`emitT` on a positive answer and `emitF` on a negative answer; in every other state it
has no call.
\end{definition}

\begin{definition}[Initial work-tape configuration]
\lean{Complexity.ImplicitEnum.cfg}
\label{Complexity.ImplicitEnum.cfg}
\leanok
\uses{Complexity.ImplicitEnum.St, Turing.Cfg, Turing.FinTM.bufferTape}
For Boolean lists $x$ and $out$, a state $s$ of the specified seven-state set, and a
natural number $i$, define a configuration on input $x$ with state $s$, input-head
position $1$, output $out$, and work-tape head at position $0$; the work tape contains
the binary expansion of $i$ as its Boolean tape contents.
\end{definition}

\begin{lemma}[Encoding a register call]
\lean{Complexity.ImplicitEnum.vword_call}
\label{Complexity.ImplicitEnum.vword_call}
\leanok
\uses{Complexity.ImplicitEnum.St, Complexity.ImplicitEnum.cfg, Complexity.LogProg.Mode.seg0, Complexity.LogProg.argSegs, Complexity.LogProg.callSegs, Complexity.LogProg.regWords, Complexity.LogProg.render, Complexity.LogProg.tapeWord_bufferTape, Complexity.LogProg.vword, Turing.pairEncode, Turing.pairEncode_eq_dbl}
\difficulty{1}
For any Boolean word $x$, output word $u$, state $s$, natural number $i$, decider index,
and positive and negative successor states, the virtual input formed by calling the
whole-input decider on the input segment $x$ and the single register containing the
binary representation of $i$ is the self-delimiting encoding $\langle
x,\operatorname{bits}(i)\rangle$.
\end{lemma}

\begin{lemma}[State of an initial configuration]
\lean{Complexity.ImplicitEnum.cfg_state}
\label{Complexity.ImplicitEnum.cfg_state}
\leanok
\uses{Complexity.ImplicitEnum.St, Complexity.ImplicitEnum.cfg}
\difficulty{1}
For any Boolean input word $x$, output word $o$, state $s$ of the specified finite state
set, and natural number $i$, the state component of the configuration constructed from
$x$, $o$, $s$, and $i$ is $s$.
\end{lemma}

\begin{lemma}[Transition from the asking state]
\lean{Complexity.ImplicitEnum.rstep_ask}
\label{Complexity.ImplicitEnum.rstep_ask}
\leanok
\uses{Complexity.ImplicitEnum.cfg, Complexity.ImplicitEnum.cfg_state, Complexity.ImplicitEnum.prog, Complexity.ImplicitEnum.vword_call, Complexity.LogProg.rstep, Turing.pairEncode}
\difficulty{3}
For any Boolean word $x$, output word $u$, natural number $i$, and Boolean oracle $o$ on
encoded Boolean words, one step of the enumerator from the configuration with output
$u$, state asking, and counter $i$ leaves the output and counter unchanged and enters
the bit state if $o(\langle x,\operatorname{bits}(i)\rangle)$ is true, and otherwise
enters the halted state.
\end{lemma}

\begin{lemma}[Transition from the bit state]
\lean{Complexity.ImplicitEnum.rstep_bit}
\label{Complexity.ImplicitEnum.rstep_bit}
\leanok
\uses{Complexity.ImplicitEnum.cfg, Complexity.ImplicitEnum.cfg_state, Complexity.ImplicitEnum.prog, Complexity.ImplicitEnum.vword_call, Complexity.LogProg.rstep, Turing.pairEncode}
\difficulty{3}
For any Boolean word $x$, output word $u$, natural number $i$, and Boolean oracle $o$ on
encoded Boolean words, one step of the enumerator from the configuration with output
$u$, state holding a bit, and counter $i$ leaves the output and counter unchanged and
enters the state that emits true if $o(1,\langle x,\operatorname{bits}(i)\rangle)$ is
true, and the state that emits false otherwise.
\end{lemma}

\begin{definition}[Shape condition for configurations]
\lean{Complexity.ImplicitEnum.Shape}
\label{Complexity.ImplicitEnum.Shape}
\leanok
\uses{Complexity.ImplicitEnum.St, Complexity.ImplicitEnum.cfg, Complexity.ImplicitEnum.prog, Turing.Cfg}
Let $f$ map Boolean lists to Boolean lists, and let $o_0,o_1$ map a Boolean list and a
natural number to a Boolean. Assume that for every Boolean list $x$ and natural number
$i$, $o_0$ on the encoding of $x$ paired with the binary expansion of $i$ is true
exactly when $i$ is less than the length of $f(x)$, and $o_1$ on that encoding equals
the $i$th entry of $f(x)$, with false used when that entry is out of range. For Boolean
list $x$, natural numbers $F,W$, and a machine configuration $c$ on input $x$, either
$c$ is the configuration defined above for some output list, state `ask` or `bit`, and
natural number $i\le F$, or every active state of $c$ has no call and the head of its
work tape lies between $-1$ and $W$, inclusive.
\end{definition}

\begin{lemma}[Register update on an enumerator configuration]
\lean{Complexity.ImplicitEnum.regCfg_cfg}
\label{Complexity.ImplicitEnum.regCfg_cfg}
\leanok
\uses{Complexity.ImplicitEnum.St, Complexity.ImplicitEnum.cfg, Complexity.LogProg.regCfg, Turing.FinTM.bufferTape}
\difficulty{4}
Let $f$ be a function from Boolean words to Boolean words, let $o$ assign a Boolean
value to each pair consisting of an index in $\{0,1\}$ and a Boolean word, and suppose
that for every Boolean word $x$ and natural number $i$, $o(0,\langle
x,\operatorname{bits}(i)\rangle)$ records whether $i$ is less than the length of $f(x)$,
while $o(1,\langle x,\operatorname{bits}(i)\rangle)$ is the $i$-th bit of $f(x)$,
defaulting to false when that position is absent. For any Boolean words $x,o'$, states
$s,s'$ of the specified finite state set, and natural numbers $i,n$, updating register
$0$ of the configuration for $x,o',s,i$ to state $s'$, tape content storing the binary
expansion of $n$, and head position $0$ yields the configuration for $x,o',s',n$.
\end{lemma}

\begin{lemma}[One output bit of the enumeration]
\lean{Complexity.ImplicitEnum.round}
\label{Complexity.ImplicitEnum.round}
\leanok
\uses{Complexity.ImplicitEnum.Shape, Complexity.ImplicitEnum.cfg, Complexity.ImplicitEnum.prog, Complexity.ImplicitEnum.regCfg_cfg, Complexity.ImplicitEnum.rstep_ask, Complexity.ImplicitEnum.rstep_bit, Complexity.ImplicitEnum.tm, Complexity.LogProg.inc_run, Complexity.LogProg.regCfg, Complexity.LogProg.rrun, Complexity.LogProg.rrun_add, Complexity.LogProg.rrun_one, Complexity.LogProg.rrun_zero, Complexity.LogProg.rstep_noncall, Turing.Action.apply, Turing.MultiTapeTM.step}
\difficulty{5}
Let $f$ map Boolean words to Boolean words, and let $o_0,o_1$ be Boolean-valued oracles
on Boolean words such that, for every Boolean word $x$ and natural number $j$,
$o_0(\langle x,\operatorname{bits}(j)\rangle)$ is true exactly when $j<|f(x)|$, and
$o_1(\langle x,\operatorname{bits}(j)\rangle)$ equals the $j$-th bit of $f(x)$,
defaulting to false when that bit is absent. If $i<|f(x)|$, $|f(x)|\le F$, and
$|\operatorname{bits}(i+1)|\le W$, then the enumerator has a finite run from the
configuration with output $f(x)$'s first $i$ bits, state asking, and counter $i$ to the
configuration with output its first $i+1$ bits, state asking, and counter $i+1$; every
configuration strictly before the end of this run has shape bounded by $F$ and $W$.
\end{lemma}

\begin{lemma}[Successive enumeration rounds]
\lean{Complexity.ImplicitEnum.rounds}
\label{Complexity.ImplicitEnum.rounds}
\leanok
\uses{Complexity.ImplicitEnum.Shape, Complexity.ImplicitEnum.cfg, Complexity.ImplicitEnum.prog, Complexity.ImplicitEnum.round, Complexity.LogProg.rrun, Complexity.LogProg.rrun_add, Complexity.LogProg.rrun_zero}
\difficulty{4}
Let $f$ and the two Boolean oracles satisfy, for every Boolean word $x$ and natural
number $j$, $o_0(\langle x,\operatorname{bits}(j)\rangle)$ is true exactly when
$j<|f(x)|$, and $o_1(\langle x,\operatorname{bits}(j)\rangle)$ equals the $j$-th bit of
$f(x)$, defaulting to false when absent. Suppose $|f(x)|\le F$, and
$|\operatorname{bits}(j+1)|\le W$ for every $j<|f(x)|$. For every $i\le |f(x)|$, the
enumerator has a finite run from its initial configuration with empty output and counter
zero to a configuration with the first $i$ bits of $f(x)$ as output, state asking, and
counter $i$; every configuration strictly before the end of this run has shape bounded
by $F$ and $W$.
\end{lemma}

\begin{lemma}[Correctness of the enumerator run]
\lean{Complexity.ImplicitEnum.run}
\label{Complexity.ImplicitEnum.run}
\leanok
\uses{Complexity.ImplicitEnum.Shape, Complexity.ImplicitEnum.St, Complexity.ImplicitEnum.cfg, Complexity.ImplicitEnum.cfg_state, Complexity.ImplicitEnum.prog, Complexity.ImplicitEnum.rounds, Complexity.ImplicitEnum.rstep_ask, Complexity.ImplicitEnum.tm, Complexity.LogProg.rrun, Complexity.LogProg.rrun_add, Complexity.LogProg.rrun_one, Complexity.LogProg.rstep_noncall, Turing.Action.apply, Turing.Cfg, Turing.Cfg.init, Turing.MultiTapeTM.step}
\difficulty{5}
Let $f$ map Boolean words to Boolean words, and let $o_0,o_1$ be Boolean-valued oracles
such that, for every Boolean word $x$ and natural number $j$, $o_0(\langle
x,\operatorname{bits}(j)\rangle)$ is true exactly when $j<|f(x)|$, and $o_1(\langle
x,\operatorname{bits}(j)\rangle)$ equals the $j$-th bit of $f(x)$, defaulting to false
when absent. If $|f(x)|\le F$ and $|\operatorname{bits}(j+1)|\le W$ for every
$j<|f(x)|$, then the enumerator, started in its initial configuration on $x$, reaches
after some finite number $N$ of steps a halted configuration whose output is $f(x)$ and
whose work-tape head is at position zero. Every configuration at a time strictly less
than $N$ has shape bounded by $F$ and $W$.
\end{lemma}

\begin{theorem}[Space-bounded computation from implicit logspace]
\lean{Complexity.ImplicitlyLogspaceComputable.computesInSpace}
\label{Complexity.ImplicitlyLogspaceComputable.computesInSpace}
\leanok
\uses{Complexity.ImplicitEnum.cfg, Complexity.ImplicitEnum.cfg_state, Complexity.ImplicitEnum.prog, Complexity.ImplicitEnum.run, Complexity.ImplicitEnum.vword_call, Complexity.ImplicitlyLogspaceComputable, Complexity.LogProg.CallOK, Complexity.LogProg.CleanRun.mono, Complexity.LogProg.bankK, Complexity.LogProg.bankStart, Complexity.LogProg.bankTM, Complexity.LogProg.bank_cleanRun, Complexity.LogProg.callSegs, Complexity.LogProg.compileFinTM, Complexity.LogProg.compile_space, Complexity.LogProg.length_bits_le_log, Complexity.LogProg.log_poly_bound, Complexity.LogProg.regWords, Complexity.LogProg.rrun, Complexity.LogProg.tapeWord_bufferTape, Complexity.LogProg.vword, Complexity.indexLang, Complexity.logSpace, Complexity.logSpace_mono, Complexity.pairEncode_mem_indexLang, Turing.Cfg.init, Turing.FinTM, Turing.FinTM.ComputesInSpace, Turing.FinTM.DecidesInSpace, Turing.MultiTapeTM.indicator, Turing.pairEncode}
\difficulty{7}
Let $f:\{0,1\}^*\to\{0,1\}^*$ have polynomially bounded output length, and
suppose that both languages of encoded pairs $(x,i)$ for which the $i$th output bit
of $f(x)$ is $1$, and for which $i<|f(x)|$, belong to $\text{LOGSPACE}$. Then there
are a finite-state deterministic Turing machine $M$ and a constant $c\in\bbn$ such
that, on every input $x$, $M$ outputs $f(x)$ and uses at most $c\log |x|$ work-tape
cells.
\end{theorem}

\begin{theorem}[Polynomial-time computability from implicit logspace]
\lean{Complexity.ImplicitlyLogspaceComputable.polyTimeComputable}
\label{Complexity.ImplicitlyLogspaceComputable.polyTimeComputable}
\leanok
\uses{Complexity.ImplicitlyLogspaceComputable, Complexity.ImplicitlyLogspaceComputable.computesInSpace, Complexity.PolyTimeComputable, Complexity.polyTimeComputable_of_computesInSpace}
\difficulty{1}
Let $f:\{0,1\}^*\to\{0,1\}^*$ have polynomially bounded output length, and
suppose that both languages of encoded pairs $(x,i)$ for which the $i$th output bit
of $f(x)$ is $1$, and for which $i<|f(x)|$, belong to $\text{LOGSPACE}$. Then there
are a finite-state deterministic Turing machine $M$ and constants $C,c\in\bbn$ such
that on every input of length $n$, $M$ outputs $f(x)$ within $C(n+1)^c$ steps.
\end{theorem}
